\documentclass[10pt,a4paper]{amsart}
\usepackage[margin=19mm]{geometry}
\usepackage{amsmath,amssymb,mathtools}
\usepackage[scr=rsfs,cal=euler]{mathalfa}
\usepackage{xfrac}
\usepackage[unicode,hidelinks,bookmarks=false]{hyperref}
\newcommand{\ie}{i.e.\ }
\newcommand{\cf}{cf.\ }
\newcommand{\resp}{respectively\ }
\newcommand{\Subsection}{\subsection}
\providecommand{\nobreakdash}{\nobreak\hbox{-}}
\setlength{\emergencystretch}{2em}
\multlinegap0pt
\usepackage{bm}
\usepackage{bbm}
\usepackage{dsfont}
\usepackage{xspace}
\usepackage{cancel}
\usepackage{mathtools}
\usepackage{amsthm}
\makeatletter
\@ifundefined{theorem}{\newtheorem{theorem}{Theorem}}{}
\@ifundefined{proposition}{\newtheorem{proposition}[theorem]{Proposition}}{}
\makeatother
\title[Smooth geometry of skew PBW extensions over commutative poly]{Smooth geometry of skew PBW extensions over commutative polynomial rings I}
\author{Andr\'es Rubiano, Armando Reyes}
\date{Source: arXiv:2409.18947v2 (CC BY 4.0)}
\begin{document}
\maketitle

\noindent\emph{Source and licence.} Smooth geometry of skew PBW extensions over commutative polynomial rings I. Version arXiv:2409.18947v2 (CC BY 4.0), distributed under the Creative Commons Attribution 4.0 International licence, as recorded in the arXiv licence metadata for this exact version and shown on the abstract page https://arxiv.org/abs/2409.18947v2. Full rights notice: licenses/arxiv-2409.18947-CC-BY-4.0.txt.\par\smallskip

\noindent\textit{Editorial scope.} Counted unit: the Proposition whose source label is \texttt{autoSkewPBW23} in the article (source0.tex lines 2050--2085, source label \texttt{autoSkewPBW23}), delivered together with the article's own complete original proof of it (source0.tex lines 2086--2223, one \texttt{proof} environment of 138 lines). No mathematics is written, repaired or re-derived: statement and proof are the article's own text line for line. Required context retained uncounted: DSdostresSPBW1 (lines 2043--2046). Every definition, display and label of those lines is retained unchanged. Omitted from this package and not redistributed: 1--2042, 2047--2049, 2224--2930. Nothing inside a retained excerpt is omitted or altered: every retained body is delivered exactly as the article's lines are written, so no omission receipt of any kind accompanies this package. Citations: the retained excerpts carry no citation command at all, so no bibliography entry of the article is transcribed and no reference is redistributed. Reference adaptation: none was needed and none was made; every reference inside the delivered unit resolves inside it, so no unresolved reference or citation is produced. Printed slips kept literal: no printed word, symbol, hypothesis or number of the counted statement or of its proof was altered. Layout and class: the article's own class is \texttt{amsart} with packages including amsthm, amssymb, amsfonts, latexsym, mathtools, thmtools, amsmath, lscape, fontenc, tikz-cd, enumitem, longtable, mathrsfs, array; the wrapper is the lane's pinned \texttt{amsart} editorial wrapper at 10pt on A4, which restates the theorem-like environments the retained text needs and adds the article's own macro definitions that the retained text needs (none), with extra wrapper packages (bm, bbm, dsfont, xspace, cancel, mathtools). The delivered numbering is the unit's own. Assets: no retained span contains a figure environment, a graphics-inclusion command, a PDF-page inclusion, a table or a TikZ picture; no image file is copied into this package, the package declares no asset dependency and no image file is redistributed with it. Licence: version arXiv:2409.18947v2 of this article is distributed under the Creative Commons Attribution 4.0 International licence, as recorded in the arXiv licence metadata for this exact version and shown on the abstract page https://arxiv.org/abs/2409.18947v2. A full rights notice is delivered at licenses/arxiv-2409.18947-CC-BY-4.0.txt. Characters: every retained excerpt is pure ASCII, so the delivered engine, XeLaTeX with its default Latin Modern text font, reports no missing character.
\begin{align}
    x_it_j = &\  a_{ijj}t_jx_i+b_{ij}x_i+ p_i, \quad \text{ for } i=1, 2, 3 \text{ and } j=1, 2, \label{DSdostresSPBW11} \\
    x_j x_i = &\ c_{i, j}x_i x_j + q_{i,j}^{(0)} + \sum_{k=1}^{3}q_{i,j}^{(k)}x_k, \quad\ {\rm for}\ i, j, k \in \{1, 2, 3\}, \ {\rm and} \ i < j, \label{DSdostresSPBW1}
\end{align}

\begin{proposition}\label{autoSkewPBW23}
Let
\begin{align}
   \nu_{t_i}(t_j) = &\ t_j, & \nu_{t_i}(x_k) = &\ a_{kii}x_k, \label{Autoskew4.38} \\
   \nu_{x_k}(t_i) = &\ a_{kii}^{-1}(t_i-b_{ki}),   &  \nu_{x_k}(x_k) = &\ x_k, \label{Autoskew4.39} \\
   \nu_{x_i}(x_j) = &\ c_{i,j}x_j+q_{i,j}^{(i)}, \quad {\rm for}\ i < j, & \nu_{x_i}(x_j) = &\ c_{j,i}^{-1}x_j-c_{j,i}^{-1}q_{j,i}^{(i)},  \quad {\rm for}\ i > j,  \label{Autoskew4.40}   
\end{align}

Then:
\begin{enumerate}
\item [\rm (1)] Leibniz's rule holds if the following relationships hold:
\begin{align}
    b_{sl}(a_{ill}-1) &\ =b_{il}(a_{sll}-1), \label{cond4.563} \\
    q_{s,i}^{(s)}(a_{ill}-1) &\ =0,\label{cond4.573} \\
    p_i(c_{s,i}-a_{sll}) &\ =b_{il}q_{s,i}^{(s)}, \label{cond4.583} \\
    q_{i,j}^{(s)}(c_{s,j}c_{s,i}-1) &\ =0, \label{cond4.593} \\
    q_{s,j}^{(s)}(c_{i,j}-1) &\ = q_{i,j}^{(i)}(c_{s,j}-1), \label{cond4.603} \\
    q_{s,i}^{(s)}(c_{i,j}-1) &\ = q_{i,j}^{(j)}(c_{s,i}-1), \label{cond4.613} \\
    q_{i,j}^{(k)}(c_{s,j}c_{s,i}-c_{k,s}^{-1}) &\ =0, \quad 1\leq k \leq 3, \quad {\rm and}  \label{cond4.623} \\
     \sum_{k=1}^{s-1}c_{k,s}^{-1}q_{i,j}^{(k)}q_{k,s}^{(s)}-\sum_{k=s+1}^{n}q_{i,j}^{(k)}q_{k,s}^{(s)} &\ 
     +q_{s,i}^{(s)}q_{s,j}^{(s)}(1-c_{i,j})+(c_{s,j}c_{s,i}-1)q_{i,j}^{(0)}  = 0. \label{cond4.633}
\end{align}

for $1 \leq i,j,s \leq n$, $i<j$, $l \in \{1, 2\}$, where $c_{r,t}:=c_{t,r}^{-1}$, $q_{r,t}^{(p)}:=-c_{t,r}^{-1}q_{t,r}^{(p)}$, for all $1 \leq r,t,p \leq n$, $t<r$, and  $c_{r,r}=1$ and $q_{r,r}^{(p)}=0$ for all $1 \leq r,p \leq n$.

The maps defined by {\rm (}\ref{Autoskew4.38}{\rm )}, {\rm (}\ref{Autoskew4.39}{\rm )} and {\rm (}\ref{Autoskew4.40}{\rm )} simultaneously extend to algebra automorphisms $\nu_{t_1}, \nu_{t_2}, \nu _{x_1}, \nu_{x_2}, \nu_{x_3}$ of $\sigma(\Bbbk[t_1, t_2])\langle x_1, x_2, x_3\rangle$ when the previous relations are satisfied.
\end{enumerate}
\begin{enumerate}
\item [\rm (2)] If the mentioned conditions in (1) hold, then we have that 
\begin{equation}\label{Eq1skew23}
   \nu_{t_i} \circ \nu_{t_j} = \nu_{t_j} \circ \nu_{t_i}, \quad \nu_{t_j} \circ \nu_{x_k} = \nu_{x_k} \circ \nu_{t_j}, \quad \nu_{x_k} \circ \nu_{x_l} = \nu_{x_l} \circ \nu_{x_k}, 
\end{equation}

for $i, j = 1, 2$ and $1\le k, l \leq 3$.
\end{enumerate}
\end{proposition}

\begin{proof}
For the first assertion, the map $\nu_{t_l}$, $l=1, 2$ can be extended to an algebra homomorphism if and only if the definitions of $\nu_{t_l}(t_j)$,  and $\nu_{t_l}(x_i)$ respect relations  {\rm (}\ref{DSdostresSPBW1}{\rm )}: 
\begin{equation*}
\nu_{t_l}(x_i)\nu_{t_l}(t_j)-\nu_{t_l}(a_{ijj}t_j+b_{ij})\nu_{t_l}(x_i) = \nu_{t_l}(p_i), \quad \text{for}\ i\in\{1, 2, 3\},\ j, l \in \{1, 2\},
\end{equation*}

and
\begin{equation*}
\nu_{t_l}(x_j)\nu_{t_l}(x_i)-c_{i,j}\nu_{t_l}(x_i)\nu_{t_l}(x_j) = q_{i,j}^{(0)}+\sum_{k=1}^{3}q_{i,j}^{(k)}\nu_{t_l}(x_k), 
\end{equation*}

for $i, j \in \{1, 2, 3\}, \ l \in \{1, 2\}$ with $i < j$. Then
\begin{align}
p_i(a_{ill}-1) = &\ 0, \quad {\rm for}\ i \in\{1, 2, 3\}, \ l \in \{1, 2\}, \notag \\
q_{i,j}^{(0)}(a_{jll}a_{ill}-1) = &\ 0, \quad {\rm and} \label{firstsecondrel33} \\
q_{i,j}^{(l)}(a_{jll}a_{ill}-a_{kll}) = &\ 0, \quad {\rm for} \ i, j,k \in \{1, 2, 3\}, \ l \in \{1, 2\}. \notag 
\end{align}

It is necessary that $\nu_{x_s}$ for $s = 1, 2, 3$ can be extended to $\sigma(\Bbbk[t_1, t_2])\langle x_1, x_2, x_3\rangle$, that is, 
\begin{equation}\label{mapvxt3}
\nu_{x_s}(x_i)\nu_{x_s}(t_j)-\nu_{x_s}(a_{ijj}t_j+b_{ij})\nu_{x_s}(x_i) = \nu_{x_s}(p_i), \quad \text{for}\ i\in\{1, 2, 3\},\ j, l \in \{1, 2\},
\end{equation}

and
\begin{equation}\label{mapvxx33}
\nu_{x_s}(x_j)\nu_{x_s}(x_i)-c_{i,j}\nu_{x_s}(x_i)\nu_{x_s}(x_j) = q_{i,j}^{(0)}+\sum_{k=1}^{3}q_{i,j}^{(k)}\nu_{x_s}(x_k).
\end{equation}

From Equation {\rm (}\ref{mapvxt3}{\rm )} we get the following three cases to be considered:
\begin{itemize}
    \item $s < i$:
    \begin{align*}
     \nu_{x_s}(x_i)\nu_{x_s}(t_j) &\ - \nu_{x_s}(a_{ijj}t_j+b_{ij})\nu_{x_s}(x_i) = \nu_{x_s}(p_i)\\
     c_{s,i}(x_i+q_{s,i}^{(s)})a_{sjj}^{-1}(t_j-b_{sj}) &\ - a_{ijj}a_{sjj}^{-1}(t_j-b_{sj})c_{s,i}(x_i+q_{s,i}^{(s)}) \\
        &\ - b_{ij}c_{s,i}(x_i+q_{s,i}^{(s)}= p_i.
    \end{align*}
    
   In this way, 
    \begin{align*}
        b_{sj}(a_{ijj}-1) &\ =b_{ij}(a_{sjj}-1), \\
        q_{s,i}^{(s)}(a_{ijj}-1) &\ =0, \quad {\rm and} \\
        p_i(c_{s,i}-a_{sjj}) &\ =b_{ij}q_{s,i}^{(s)}.
    \end{align*}
    
    \item $s = i$: It is straightforward that no new conditions are obtained.
    
    \item $s > i$: 

    \begin{align*}
          \nu_{x_s}(x_i)\nu_{x_s}(t_j) &\ - \nu_{x_s}(a_{ijj}t_j+b_{ij})\nu_{x_s}(x_i) = \nu_{x_s}(p_i) \\
           c_{i,s}^{-1}(x_i-q_{i,s}^{(s)})a_{sjj}^{-1}(t_j-b_{sj}) &\ -a_{ijj}a_{sjj}^{-1}(t_j-b_{sj})c_{i,s}^{-1}(x_i-q_{i,s}^{(s)}) \\
           &\ - b_{ij}c_{i,s}^{-1}(x_i-q_{i,s}^{(s)}) = p_i.
    \end{align*}
  
Then, 
\[
q_{i,s}^{(s)}(a_{ijj}-1) = 0 \quad {\rm and} \quad p_i(c_{i,s}a_{sjj}-1) = a_{sjj}b_{ij}q_{i,s}^{(s)}. 
\]
\end{itemize}

Due to expression {\rm (}\ref{mapvxx33}{\rm )} we have to consider the following cases:
\begin{itemize}
    \item $s < i$:
    \begin{equation*}
        \nu_{x_s}(x_j)\nu_{x_s}(x_i)-c_{i,j}\nu_{x_s}(x_i)\nu_{x_s}(x_j) = q_{i,j}^{(0)}+\sum_{k=1}^{3}q_{i,j}^{(k)}\nu_{x_s}(x_k)
    \end{equation*}
    From this equation, we obtain the following new conditions
    \begin{align*}
        q_{i,j}^{(s)}(c_{s,j}c_{s,i}-1) &\ =0, \\
        q_{s,j}^{(s)}(c_{i,j}-1) &\ = q_{i,j}^{i}(c_{s,j}-1), \\
        q_{s,i}^{(s)}(c_{i,j}-1) &\ = q_{i,j}^{j}(c_{s,i}-1), \\
        q_{i,j}^{(k)}(c_{s,j}c_{s,i}-c_{k,s}^{-1}) &\ =0, \quad 1\leq k \leq s-1, \\
        q_{i,j}^{(k)}(c_{s,j}c_{s,i}-c_{s,k}) &\ =0, \quad s+1\leq k \leq 3, k\ne i, j, \quad {\rm and}  \\
        \sum_{k=1}^{s-1}c_{k,s}^{-1}q_{i,j}^{(k)}q_{k,s}^{(s)}-\sum_{k=s+1}^{3}q_{i,j}^{(k)}q_{k,s}^{(s)} &\ +q_{s,i}^{(s)}q_{s,j}^{(s)}(1-c_{i,j})+(c_{s,j}c_{s,i}-1)q_{i,j}^{(0)}  = 0.
    \end{align*}
    
     \item $i < s < j$: 
    \begin{equation*}
        \nu_{x_s}(x_j)\nu_{x_s}(x_i)-c_{i,j}\nu_{x_s}(x_i)\nu_{x_s}(x_j) = q_{i,j}^{(0)}+\sum_{k=1}^{3}q_{i,j}^{(k)}\nu_{x_s}(x_k).
    \end{equation*}
    
    Hence, 
    \begin{align*}
        q_{i,j}^{(s)}(c_{s,j}-c_{i,s}) &\ =0, \\
        q_{s,j}^{(s)}(c_{i,j}-1) &\ = q_{i,j}^{(i)}(c_{s,j}-1), \\
        q_{i,s}^{(s)}(c_{i,j}-1) &\ = q_{i,j}^{(j)}(c_{i,s}-1), \\
        q_{i,j}^{(k)}(c_{s,j}c_{i,s}^{-1}-c_{k,s}^{-1}) &\ =0, \quad 1\leq k \leq s-1, k \ne i,  \\
        q_{i,j}^{(k)}(c_{s,j}c_{i,s}^{-1}-c_{s,k}) &\ =0, \quad s+1\leq k \leq n, k\ne j, \quad {\rm and} \\
        \sum_{k=1}^{s-1}c_{k,s}^{-1}q_{i,j}^{(k)}q_{k,s}^{(s)} -  \sum_{k=s+1}^{3}q_{i,j}^{(k)}q_{s,k}^{(s)} &\ + c_{i,s}^{-1}q_{s,j}^{(s)}q_{i,s}^{(s)}(c_{i,j}-1)  +(c_{s,j}c_{i,s}^{-1}-1)q_{i,j}^{(0)}  = 0.
    \end{align*}
    
    \item $s\geq j$: 
    \begin{equation*}
        \nu_{x_s}(x_j)\nu_{x_s}(x_i)-c_{i,j}\nu_{x_s}(x_i)\nu_{x_s}(x_j) = q_{i,j}^{(0)}+\sum_{k=1}^{3}q_{i,j}^{(k)}\nu_{x_s}(x_k).
    \end{equation*}

Then, 
    \begin{align*}
        q_{i,j}^{(s)}(c_{j,s}c_{i,s}-1) &\ =0, \\
        q_{j,s}^{(s)}(c_{i,j}-1) &\ = q_{i,j}^{(i)}(c_{j,s}-1), \\
        q_{i,j}^{(k)}(c_{j,s}^{-1}c_{i,s}^{-1}-c_{k,s}^{-1}) &\ =0, \quad 1\leq k \leq s-1, k \ne i, j\\
        q_{i,j}^{(k)}(c_{j,s}^{-1}c_{i,s}^{-1}-c_{s,k}) &\ =0, \quad s+1\leq k \leq n,  \quad {\rm and} \\
        \sum_{k=1}^{s-1}c_{k,s}^{-1}q_{i,j}^{(k)}q_{k,s}^{(s)} - \sum_{k=s+1}^{3}q_{i,j}^{(k)}q_{s,k}^{(s)}  &\ + c_{j,s}^{-1}c_{i,s}^{-1}q_{j,s}^{(s)}q_{i,s}^{(s)}(1-c_{i,j})  +(c_{j,s}^{-1}c_{i,s}^{-1}-1)q_{i,j}^{(0)}  = 0.
    \end{align*}
\end{itemize}

If we put together all the conditions for $s \in \{1, 2, 3\}$, then we obtain the restrictions for the extension of the automorphisms.

For the second assertion, it is enough to prove it for the generators $t_j$, $x_i$, $1\leq i \leq 3$ and $j \in \{1, 2\}$. Since
\begin{align}
    \nu_{t_k} \circ \nu_{t_m}(t_j) = &\ \nu_{t_k}(t_j)=t_j, \label{comptttj3}\\
    \nu_{t_m} \circ \nu_{t_k}(t_j) = &\ \nu_{t_m}(t_j)=t_j, \label{comptttj3.} \\
    \nu_{t_k} \circ \nu_{t_m}(x_i) = &\ \nu_{t_k}(a_{imm}x_i)=a_{imm}a_{ikk}x_i, \quad {\rm and}  \label{compttxi3}\\
    \nu_{t_m} \circ \nu_{t_k}(x_i) = &\ \nu_{t_m}(a_{ikk}x_i)=a_{imm}a_{ikk}x_i. \label{compttxi3.}
\end{align}

for all $1\leq i \leq 3$ and $k,m,j \in \{1, 2\}$, then all relations are satisfied, and hence $\nu_{t_k} \circ \nu_{t_m} = \nu_{t_m} \circ \nu_{t_k}$ for $k,m \in \{1, 2\}$.

Now, 
\begin{align}
    \nu_{t_k} \circ \nu_{x_m}(t_j) = &\ \nu_{t_k}(a_{mjj}^{-1}(t_j-b_{mj}))=a_{mjj}^{-1}(t_j-b_{mj}), \label{comptxtj3}\\
    \nu_{x_m} \circ \nu_{t_k}(t_j) = &\ \nu_{x_m}(t_j)=a_{mjj}^{-1}(t_j-b_{mj}), \label{comptxtj3.} \\
    \nu_{t_k} \circ \nu_{x_m}(x_i) = &\ \nu_{t_k}(c_{m,i}x_i+q_{m,i}^{(m)})=c_{m,i}a_{ikk}x_i+q_{m,i}^{(m)}, \label{comptxxi3} \quad {\rm and} \\
    \nu_{x_m} \circ \nu_{t_k}(x_i) = &\ \nu_{x_m}(a_{ikk}x_i)=a_{ikk}(c_{m,i}x_i+q_{m,i}^{(m)}), \label{comptxxi3.}
\end{align}

for all $1\leq i,m \leq 3$ and $k,j \in \{1, 2\}$. As it is clear, expressions {\rm (}\ref{comptxtj3}{\rm )} and {\rm (}\ref{comptxtj3.}{\rm )} hold, while relations {\rm (}\ref{comptxxi3}{\rm )} and {\rm (}\ref{comptxxi3.}{\rm )} are also satisfied due to expression {\rm (}\ref{cond4.593}{\rm )}. In this way, $\nu_{t_k} \circ \nu_{x_m} = \nu_{x_m} \circ \nu_{t_k}$, for  $1\leq m \leq 3$ and $k \in \{1, 2\}$.

Finally, 
\begin{align}
    \nu_{x_k} \circ \nu_{x_m}(t_j) = &\ \nu_{x_k}(a_{mjj}^{-1}(t_j-b_{mj}))=a_{mjj}^{-1}a_{kjj}^{-1}(t_j-b_{kj})-a_{mjj}^{-1}b_{jm}, \label{compxxtj3}\\
    \nu_{x_m} \circ \nu_{x_k}(t_j) = &\ \nu_{x_m}(a_{kjj}^{-1}(t_j-b_{kj}))=a_{kjj}^{-1}a_{mjj}^{-1}(t_j-b_{mj})-a_{kjj}^{-1}b_{kj}, \label{compxxtj3.} \\
    \nu_{x_k} \circ \nu_{x_m}(x_i) = &\ \nu_{x_k}(c_{m,i}x_i+q_{m,i}^{(m)})=c_{m,i}(c_{k,i}x_i+q_{k,i}^{(k)})+q_{m,i}^{(m)}, \quad {\rm and} \label{compxxxi3}\\
    \nu_{x_m} \circ \nu_{x_k}(x_i) = &\ \nu_{x_m}(c_{k,i}x_i+q_{k,i}^{(k)})=c_{k,i}(c_{m,i}x_i+q_{m,i}^{(m)})+q_{k,i}^{(k)}, \label{compxxxi3.}
\end{align}

for all $1\leq i,m,k  \leq 3$ and $j \in \{1, 2\}$. Then relations {\rm (}\ref{compxxtj3}{\rm )} and {\rm (}\ref{compxxtj3.}{\rm )} hold because relation {\rm (}\ref{cond4.563}{\rm )} is satisfied. With respect to relations {\rm (}\ref{compxxtj3}{\rm )} and {\rm (}\ref{compxxtj3.}{\rm )}, both are satisfied due to the expressions {\rm (}\ref{cond4.603}{\rm )} and {\rm (}\ref{cond4.613}{\rm )}. This yields that $\nu_{x_k} \circ \nu_{x_m} = \nu_{x_m} \circ \nu_{x_k}$, for  $1\leq k, m \leq 3$.
\end{proof}

\end{document}
